\chapter{Docker: Image, Container, Compose}
\label{ch:docker}

\section{Container thực chất là gì}

Container không phải là một máy ảo thu nhỏ. Nó là một tiến trình Linux
bình thường, được cô lập bởi hai tính năng của kernel: \emph{namespace}
(góc nhìn riêng về tiến trình, mạng, các mount của hệ thống file,
hostname) và \emph{cgroups} (giới hạn CPU và bộ nhớ). Không có hệ điều
hành khách, không có hypervisor --- đó là lý do container khởi động
trong vài mili giây, và cũng là lý do JVM bên trong phải được báo về
giới hạn bộ nhớ của nó (\code{MaxRAMPercentage} của Chương~17) --- vì
nhìn từ bên trong, nó thấy RAM của \emph{máy chủ}.

\emph{Image} là hệ thống file đã đóng gói cộng với metadata (entrypoint,
biến môi trường, cổng khai báo), được build thành một chồng các
\emph{layer} bất biến. Layer được định danh theo nội dung và dùng chung:
năm mươi image \code{FROM eclipse-temurin:25-jre-alpine} chỉ lưu layer
JRE một lần. \emph{Registry} (Docker Hub, GHCR, \code{registry:2} của
dự án) lưu trữ image; \code{docker push}/\code{pull} chỉ chuyển những
layer mà phía bên kia còn thiếu.

\section{Dockerfile của dự án, có chú giải}

Service nào cũng mang đúng năm dòng này:

\begin{lstlisting}[language=dockerfile]
FROM eclipse-temurin:25-jre-alpine    (*@\co{1}@*)
WORKDIR /app
COPY target/*.jar app.jar             (*@\co{2}@*)
EXPOSE 8082                           (*@\co{3}@*)
ENTRYPOINT ["java", "-jar", "app.jar"](*@\co{4}@*)
\end{lstlisting}

\begin{description}
  \item[\co{1}] Image gốc: một JRE (không phải JDK --- lúc chạy không
  cần trình biên dịch) trên Alpine Linux. Image gốc nhỏ hơn = bề mặt tấn
  công nhỏ hơn và pull nhanh hơn.
  \item[\co{2}] Sao chép fat jar do Maven build ra. Điểm tinh tế:
  Dockerfile này đòi hỏi \code{mvn package} phải chạy \emph{trước} ---
  việc build diễn ra bên ngoài. Hệ quả: một lệnh \code{COPY} cho jar cỡ
  \(\sim\)60\,MB tạo ra một layer to, nên \emph{mọi} lần build đều push
  lại toàn bộ dù chỉ sửa một dòng. Jib (Chương~13) sinh ra để sửa đúng
  chuyện này.
  \item[\co{3}] Chỉ là tài liệu, không phải hành động --- nó không mở gì
  cả. Việc công bố cổng diễn ra lúc chạy.
  \item[\co{4}] Entrypoint dạng exec: java là PID~1, nhận tín hiệu trực
  tiếp, nên \code{docker stop} kích hoạt quá trình tắt Spring sạch sẽ.
\end{description}

\section{Mạng và Compose}

Các container trên một \emph{bridge network} của Docker tìm thấy nhau
bằng \emph{tên service} --- Docker chạy một DNS server trả lời cho
\code{postgres-course}, \code{kafka}, \code{minio}. Còn \code{localhost},
bên trong một container, chính là \emph{container đó} --- cái bẫy số
một với người mới, và là lý do mọi URL mặc định cho dev trong các file
cấu hình đều cần ghi đè theo từng môi trường. Cổng chỉ với tới được từ
máy chủ khi đã được \emph{công bố} (\code{-p 5434:5432}).

Compose biến một loạt lệnh \code{docker run} thành một file khai báo
duy nhất --- \code{docker-compose.yml} của dự án là cả stack dev: bốn
kho dữ liệu, Kafka, Zipkin, Maildev và bảy service, với ba thành ngữ
đáng gọi tên:

\begin{lstlisting}[language=yaml]
  discovery-server:
    depends_on:
      config-server:
        condition: service_healthy    (*@\co{1}@*)
  config-server:
    healthcheck:
      test: ["CMD", "curl", "-f", "http://localhost:8888/actuator/health"] (*@\co{2}@*)
  postgres-course:
    ports:
      - "5434:5432"                   (*@\co{3}@*)
    volumes:
      - postgres-course-data:/var/lib/postgresql/data
\end{lstlisting}

\co{1}+\co{2} mã hóa thứ tự khởi động --- và không chỉ ``đã chạy'' mà
là ``đã trả lời endpoint sức khỏe,'' điều làm cho thứ tự này có ý nghĩa
thật. \co{3} ánh xạ hai cổng máy chủ khác nhau (5433, 5434) vào cùng
cổng nội bộ 5432 của hai container Postgres --- tên thì riêng theo từng
mạng, cổng máy chủ thì toàn cục. Volume có tên giữ dữ liệu qua những
lần tạo lại container; Chương~19 là phiên bản Kubernetes của cùng nhu
cầu này.

\begin{decision}[vì sao Compose vẫn sống cạnh Kubernetes]
Cụm (cluster) là đích triển khai, nhưng Compose vẫn là vòng lặp bên
trong: một lệnh, chạy ngay trên laptop, cổng debug đã map sẵn (mọi
service đều mở một cổng JDWP trong Compose --- gắn debugger của IDE vào
5082 và bước từng dòng qua course-service). Kubernetes mua được tự phục
hồi và triển khai cuốn chiếu bằng cái giá là nghi thức rườm rà; một lập
trình viên đang lặp đi lặp lại trên logic nghiệp vụ không muốn trả cả
hai. Hai môi trường, một hợp đồng image chung giữa chúng.
\end{decision}

\begin{pitfall}[chạy tốt trong Compose, chết trong Kubernetes]
Triệu chứng: stack xanh ở máy local; cùng image đó lại crash-loop trong
cụm. Những nghi phạm quen thuộc, đều thuộc về môi trường: một URL vẫn
trỏ tới tên service của Compose trong khi tên Service bên k8s khác đi;
một \code{depends\_on} mà cụm không tôn trọng (Kubernetes không có thứ
tự khởi động --- Chương~15); hoặc giới hạn bộ nhớ có trong cụm nhưng
không có ở local, biến một JVM béo thành OOMKill. Image là y hệt --- hãy
so \emph{môi trường}, không phải code: \code{docker inspect} với
\code{kubectl describe pod}.
\end{pitfall}

\begin{quiz}
\begin{enumerate}
  \item Container và VM đều cô lập. Nêu hai cơ chế container dùng, và
  hai hệ quả của việc không có kernel khách.
  \item Vì sao \code{localhost:5432} hỏng bên trong container dù chạy
  được trên laptop của bạn, và Compose với Kubernetes đưa ra hai cách
  sửa khác nhau nào?
  \item Cả hai container Postgres đều lắng nghe cổng 5432 nội bộ mà
  không xung đột --- giải thích. Xung đột lại có thể xảy ra ở đâu?
  \item Điều gì khiến \code{condition: service\_healthy} mạnh hơn
  \code{depends\_on} thường, và file nào của dự án cung cấp tín hiệu
  sức khỏe mà nó chờ?
\end{enumerate}
\end{quiz}
